\section{Discussion}
\label{sec:discussion}

\paragraph{What the results show.} The findings above are descriptive: agents report outcomes and do not report the counterfactuals their instructions foreclosed. Whether that is a defect is an interpretation, and we set out the case for reading it as one.

\paragraph{Omission without a motive.} The agents here were loyal in the sense most evaluations measure. They followed the client's requirement, protected the confidential figure, reached deals above the client's fallback, and mostly stated the outcome correctly. They did not treat the client's instruction as something the client might want to revisit. A requirement is a guess about what the principal values, and the negotiation prices that guess; a report that says only whether the requirement was met withholds the price.

This separates the result from the deception literature, where the agent has a reason to withhold: a failure to hide \citep{advani2026falsesuccess,guo2025upward}, misconduct to conceal \citep{scheurer2023strategic}, or a sponsor to serve \citep{wu2026adschatbots}. Here disclosure would, if anything, excuse a weaker outcome. In cheap-talk terms \citep{crawford1982strategic}, sender and receiver preferences are aligned and the message is costless, which is the configuration in which fully revealing communication is sustainable, and LLM senders are known to over-reveal rather than coarsen when interests conflict \citep{hasanibalyani2026truthful}. The agents nonetheless transmit almost nothing about the trade-off. We read this as a default about what a report contains, not as strategic distortion.

\paragraph{Accuracy belongs to three models; disclosure belongs to developers.} Once the extractor fault described in \cref{app:deviations} is removed, stated-value error is a property of three models rather than a gradient: every frontier model from every developer states its deal's value correctly in nearly every report, and the errors sit with Claude Haiku 4.5, DeepSeek V4 Pro and Mistral Large 3. Disclosure of the blocked deal does not sort the same way. The Anthropic frontier models convey it in roughly a quarter of exposed episodes; GPT-6 Astra, an equally capable negotiator at the top of its own developer's ladder, almost never does, and neither do the largest Qwen models. Negotiated value cannot separate these models: the scripted counterparty caps what any agent can win, so every model lands near the same gain fraction by construction, and the weak correlation between disclosure and negotiated value in this grid is a consequence of that design, not evidence about skill. Throughout, ``capability'' means a developer's own tier ordering, not measured negotiating ability, and within Anthropic that ordering also spans a generation. Two models from one developer converging on the same disclosure rate, with no other model near it, is consistent with a post-training difference between developers, but we did not manipulate training and cannot attribute the difference to it. What the data support is narrower: reporting fidelity does not arrive reliably with capability, and in this grid it varies more between developers than between tiers.

\paragraph{What the interventions do and do not show.} Two of our instruments name the thing they are looking for. The knowledge probe asks about better offers that were turned down, and the client question in the intervention arm asks about turned-down terms; a model with the transcript still in context will answer both with a little arithmetic. So those results establish that the information was recoverable and that a specific question elicits it, which rules out the explanation that the agents could not identify the blocked deal, but they do not show a change in what the agent judges worth telling. The deployer-side instruction is the better test of that, because it names no term: it tells the agent that the client decides from its account whether to keep their instructions, and leaves the agent to work out what follows. Under it, \intNormBelowHalf\ of the \nInterventionModels\ models begin to volunteer the trade-off in a majority of episodes, \intNormHalfToEighty\ move partway, and \intNormAboveEighty\ are nearly unmoved. A general norm about what a report is for is enough for some models and not for others, and nothing in a model's tier predicts which.

\paragraph{Oversight has to ask the question.} For anyone deploying a delegated agent today, the practical reading is that a specific question works for \intTradeoffsBelowHalf\ of the \nInterventionModels\ models we tested and a general instruction for fewer than half. Telling the agent its transcript will be read makes it write more and disclose no more. Oversight of delegated agents should therefore be built as structured elicitation, in which the interface poses the questions, rather than as an instruction to the agent to be forthcoming, since a principal cannot be expected to know which questions to ask.

\paragraph{Why principals may not notice, and what this paper cannot settle.} In a real delegated market, people represented by a weaker agent obtained worse prices but rated their deals as equally fair \citep{anthropic2026projectdeal}. Our results suggest a mechanism: reports name the blocked offer in \doneNamedAll\ of exposed episodes and convey its significance in \doneToldAll, so a principal reading one has little way to learn what the agent gave up. Whether principals want that information is the premise of the paper, and it is not something model behavior can establish. The simulated principal signed off on the report where the facts warranted intervening in \hFiveMissed\ of episodes, and most of its remaining disagreement with the facts was caution on a vaguer account; a skeptical reader can take that as evidence that the omission rarely matters. We take it as evidence about one model, and we leave the premise as an argument: a requirement is a guess about what the principal values, and a principal who cannot learn what the guess cost cannot revise it.
